\clearpage
\begin{explanationbox}
\textbf{第一步：先认熟悉的相似三角形，不引入双曲线。}
\begin{center}\begin{tikzpicture}[font=\small,>=stealth]
\draw[gray!45,rounded corners=3pt] (0,0) rectangle (14.2,9.5);
\node at (7.1,8.85) {先不画双曲线：同一直线上的长度比，等于水平投影比};

\begin{scope}[shift={(1.7,2.2)},scale=.85]
\coordinate (A) at (0,0);\coordinate (B) at (3,4);\coordinate (C) at (4.5,6);
\fill[blue!12] (A)--(3,0)--(B)--cycle;
\fill[orange!18] (B)--(4.5,4)--(C)--cycle;
\draw[blue!75!black,thick] (A)--(3,0)--(B)--cycle;
\draw[orange!85!black,thick] (B)--(4.5,4)--(C)--cycle;
\draw[gray] (2.8,0)--(2.8,.2)--(3,.2) (4.3,4)--(4.3,4.2)--(4.5,4.2);
\fill (A) circle(2.8pt) node[left=5pt] {$A$};
\fill (B) circle(2.8pt) node[above left=5pt] {$B$};
\fill (C) circle(2.8pt) node[right=5pt] {$C$};
\node[blue!75!black,left=5pt] at (1.5,2) {$5$};
\node[orange!85!black,left=5pt] at (3.75,5) {$5/2$};
\node[blue!75!black,below=5pt] at (1.5,0) {$3$};
\node[orange!85!black,below=5pt] at (3.75,4) {$3/2$};
\end{scope}
\node[anchor=west,align=left] at (8,6.5) {两个直角三角形相似\\[14pt]$\displaystyle\frac{AB}{BC}=\frac5{5/2}=2$\\[14pt]$\displaystyle\frac{\text{横向差}}{\text{横向差}}=\frac3{3/2}=2$};
\node[anchor=west,align=left] at (8,3.1) {同一方向、同一放大倍数\\[10pt]取比时，公共倍数约掉。};
\node at (7.1,.65) {接下来把蓝、橙两段对应为 $FM,MN$；原理不变。};

\end{tikzpicture}
\end{center}
同一直线上的两段倾斜程度相同，实际长度都是横向差乘同一个数。
因此长度比只需比较水平投影；这就是后面取比的初中几何基础。
\end{explanationbox}
\clearpage
\begin{explanationbox}
\textbf{第二步：把熟悉的分段结构放回第一页的图。}

第一页中，$F$ 是焦点，$M$ 在曲线上，$N$ 在另一条渐近线上；
蓝色是 $FM$，橙色是 $MN$。只需确认“过焦点、平行一条、交另一条”。

令 $k=b/a$，过右焦点的直线为 $y=k(x-c)$。
它让 $x-ay/b=c$ 保持不变，因此
\[
 (x-ay/b)(x+ay/b)=a^2
 \quad\Longrightarrow\quad x_M+ay_M/b=a^2/c.
\]
\textbf{平方差中一个因式固定，求交就降为一次。}
方法一把它画成平行等值线；方法二则从水平位移 $t$ 入手。
两法最终都回到上一页熟悉的共线分段比。

\medskip
\textbf{第三步：定长不等于位置固定。}
\begin{center}\begin{tikzpicture}[font=\small,>=stealth]
\draw[gray!45,rounded corners=3pt] (0,0) rectangle (14.2,10);
\node at (7.1,9.35) {固定 a，改变 b：橙色 MN 的长度不变，位置和方向会变};

\node at (7.1,8.35) {三图坐标尺度相同；取 $a=3$。};

\begin{scope}[shift={(0.6,5.5)},scale=.7]
\draw[gray,->] (0,0)--(5.7,0);
\draw[teal!75!black,thick,samples=70,domain=-3.8:1,variable=\v] plot ({3*sqrt(1+\v*\v/5.0625)},{\v});
\draw[violet,dashed] (0,0)--(3.5,-2.625);
\draw[blue!75!black,very thick] (3.75,0)--(3.075,-0.50625);
\draw[orange!85!black,very thick] (3.075,-0.50625)--(1.875,-1.40625);
\fill (3.75,0) circle(2.8pt) node[above] {$F$};
\fill (3.075,-0.50625) circle(2.8pt) node[right=4pt] {$M$};
\fill (1.875,-1.40625) circle(2.8pt) node[below left=4pt] {$N$};
\end{scope}
\node at (2.5,7.25) {$b=2.25$};
\node[orange!85!black] at (2.5,1.65) {$MN=3/2$};

\begin{scope}[shift={(5.1,5.5)},scale=.7]
\draw[gray,->] (0,0)--(5.7,0);
\draw[teal!75!black,thick,samples=70,domain=-3.8:1,variable=\v] plot ({3*sqrt(1+\v*\v/9)},{\v});
\draw[violet,dashed] (0,0)--(3.5,-3.5);
\draw[blue!75!black,very thick] (4.242640687119285,0)--(3.181980515339464,-1.0606601717798214);
\draw[orange!85!black,very thick] (3.181980515339464,-1.0606601717798214)--(2.1213203435596424,-2.1213203435596424);
\fill (4.242640687119285,0) circle(2.8pt) node[above] {$F$};
\fill (3.181980515339464,-1.0606601717798214) circle(2.8pt) node[right=4pt] {$M$};
\fill (2.1213203435596424,-2.1213203435596424) circle(2.8pt) node[below left=4pt] {$N$};
\end{scope}
\node at (7.0,7.25) {$b=3$};
\node[orange!85!black] at (7.0,1.65) {$MN=3/2$};

\begin{scope}[shift={(9.6,5.5)},scale=.7]
\draw[gray,->] (0,0)--(5.7,0);
\draw[teal!75!black,thick,samples=70,domain=-3.8:1,variable=\v] plot ({3*sqrt(1+\v*\v/16)},{\v});
\draw[violet,dashed] (0,0)--(3.5,-4.666666666666667);
\draw[blue!75!black,very thick] (5.0,0)--(3.4,-2.1333333333333333);
\draw[orange!85!black,very thick] (3.4,-2.1333333333333333)--(2.5,-3.3333333333333335);
\fill (5.0,0) circle(2.8pt) node[above] {$F$};
\fill (3.4,-2.1333333333333333) circle(2.8pt) node[right=4pt] {$M$};
\fill (2.5,-3.3333333333333335) circle(2.8pt) node[below left=4pt] {$N$};
\end{scope}
\node at (11.5,7.25) {$b=4$};
\node[orange!85!black] at (11.5,1.65) {$MN=3/2$};
\node at (7.1,.65) {定长不等于固定在原处；橙段随构型旋转、平移。};
\end{tikzpicture}
\end{center}
图中固定 $a=3$，三条橙段都长 $3/2$，却没有固定在同一位置。
一般地 $MN=a/2$；变化的是斜率、焦点和交点坐标。
\end{explanationbox}
